% 第 4 章 §4.4.4 尾（印刷页 214；PDF p.224）：Eichler–Shimura 同构的收尾
% 对应 工程文件/work/en/ch04.txt 第 3205–3255 行（文献表起于第 3259 行，另见 bib04.tex）。
% 本块编号公式：(4.65)、(4.66)；无定理类环境；无脚注；无插图。
% 记号：\mathcal{K} 为典范丛的幂（底本作花体 K），\mathbf{V}_q 为局部系统（底本作粗体 V），
%       V_q 为 SL_2(C) 的表示（斜体），\mathfrak{M} 为模空间（底本作 fraktur M）。

由于 $q\ge 3/2$，有 $H^0(X,\mathcal{K}^{1-q})\simeq H^1(X,\mathcal{K}^q)^*=\{0\}$。这意味着 $H^0(X,\mathbf{V}_q)=H^2(X,\mathbf{V}_q)=\{0\}$，而与 (4.64) 相伴的长正合列成为
\[
0\longrightarrow H^0(X,\mathcal{K}^q)\xrightarrow{\ \delta\ }H^1(X,\mathbf{V}_q)\longrightarrow H^1(X,\mathcal{K}^{1-q})\longrightarrow 0.
\]

上边缘映射 $\delta$ 称为「Eichler 积分」（Eichler integration）。这个名称的来由如下：若 $\omega$ 是 $\mathcal{K}^q$ 的整体全纯截面，则在足够小的开集 $U_i$ 上可以解非齐次方程 $Dy_i=\restr{\omega}{U_i}$。若令 $v_{i,j}=y_i-y_j$，则 $\{v_{i,j}\}$ 是取值于 $\mathbf{V}_q$ 的 1-上闭链（1-cocycle），代表 $\delta\omega$。

总之，由此得到同构（参看 [22, 29, 55]）
\begin{equation}
H^1(X,\mathbf{V}_q)\simeq H^0(X,\mathcal{K}^q)\oplus\bigl(H^0(X,\mathcal{K}^q)\bigr)^*.
\tag{4.65}
\end{equation}

等式 (4.65) 可用来描述 Betti 模空间在 $[\rho^{(n)}]$ 处的切空间（这是 Bill Goldman [25] 向我说明的）。由 Weil 对切空间的描述，
\begin{equation}
T_{[\rho^{(n)}]}\mathfrak{M}_B^{(n)}\simeq H^1(X,\operatorname{End}\mathbf{V}_q).
\tag{4.66}
\end{equation}

现在，$\mathrm{SL}_2(\C)$ 的表示都是自对偶（self-dual）的：$\mathbf{V}_q^*\simeq\mathbf{V}_q$。由张量积表示分解的 Clebsch--Gordon 规则，有
\[
\operatorname{End}\mathbf{V}_q=(\mathbf{V}_q\otimes\mathbf{V}_q^*)_{\mathrm{tr}=0}\simeq(\mathbf{V}_q\otimes\mathbf{V}_q)_{\mathrm{tr}=0}=\bigoplus_{\substack{j=2\\ j\in\Z}}^{2q-1}V_j
\]
（注意：平凡表示 $V_{3/2}$ 被无迹条件排除）。

\begin{annote}
底本此句作「the trivial representation $V_{3/2}$ is eliminated by the traceless condition」。按本节约定，$V_q$ 是 $2q-1$ 维的不可约表示，故平凡表示应为 $V_1$；此处疑为笔误。译文照排底本，不改动下标（参见疑点清单 Q074）。
\end{annote}

这个分解可以转译为局部系统的相应分解。由此得到
\[
H^1(X,\operatorname{End}\mathbf{V}_q)=\bigoplus_{\substack{j=2\\ j\in\Z}}^{2q-1}H^1(X,\mathbf{V}_j).
\]

把它与等式 (4.65)、(4.66) 结合起来，便得到
\[
T_{[\rho^{(n)}]}\mathfrak{M}_B^{(n)}\simeq\bigoplus_{j=2}^{n}H^0(X,\mathcal{K}^j)\oplus\bigl(H^0(X,\mathcal{K}^j)\bigr)^*.
\]

这应当与 (4.14) 作比较！

\begin{annote}
底本以感叹号提示回看 (4.14)（印刷页 177）：该式给出 Dolbeault 模空间在 $[E_F,\Phi_F]$ 处的切空间 $\bigoplus_{j=2}^{n}H^0(X,\mathcal{K}^j)\oplus\bigl(H^0(X,\mathcal{K}^j)\bigr)^*$，与本段末式形状相同；本段末式则是 Betti 模空间在 $[\rho^{(n)}]$ 处的切空间，两相对照正是 §4.4.4 的主旨（参见疑点清单 Q084）。
\end{annote}
